
% Preliminares en n{\'u}meros romanos (Resumen = p{\'a}gina I); la car{\'a}tula no se numera.
\pagenumbering{Roman}
\hypersetup{pageanchor=true}

% ================= RESUMEN =================
\section*{Resumen}
\addcontentsline{toc}{section}{Resumen}
\singlespacing\noindent
En la regi{\'o}n Cusco, m{\'a}s del 55\,\% de la poblaci{\'o}n tiene el quechua como lengua
materna y cerca de 70\,000 asegurados de EsSalud son quechuahablantes \parencite{essalud2023};
esta barrera idiom{\'a}tica afecta la calidad de la atenci{\'o}n en salud \parencite{begazo2026}.

El presente trabajo propone desarrollar y evaluar un sistema bidireccional de traducci{\'o}n de
voz Quechua--Espa{\~n}ol para consultas m{\'e}dicas, que encadena reconocimiento de voz (MMS),
traducci{\'o}n autom{\'a}tica (NLLB adaptado al quechua de Cusco, \sig{quz}) y s{\'i}ntesis de
voz (MMS-TTS) mediante el ajuste (\emph{fine-tuning}) de modelos preentrenados de c{\'o}digo
abierto.

La salida por voz en ambos sentidos busca incluir a la poblaci{\'o}n quechuahablante con
baja o nula alfabetizaci{\'o}n. La validaci{\'o}n se realizar{\'a} con consultas simuladas, sin datos
reales de pacientes, midiendo precisi{\'o}n (WER, BLEU, chrF), calidad de voz (MOS) y percepci{\'o}n
de utilidad. Se espera aportar el primer corpus cl{\'i}nico Quechua--Espa{\~n}ol y evidencia
emp{\'i}rica sobre la viabilidad de estas tecnolog{\'i}as en el dominio sanitario andino.

\vspace{0.3cm}
\textit{Palabras clave:} \nuevo{Traducci{\'o}n autom{\'a}tica, Quechua, Procesamiento del
lenguaje natural, Comunicaci{\'o}n m{\'e}dico-paciente, Inclusi{\'o}n digital.}
\onehalfspacing

\newpage
\section*{Abstract}
\addcontentsline{toc}{section}{Abstract}
\singlespacing\noindent
In the Cusco region, more than 55\% of the population has Quechua as their mother tongue and
nearly 70{,}000 EsSalud insured members are Quechua speakers \parencite{essalud2023}; this language
barrier affects the quality of health care \parencite{begazo2026}.

This work proposes to develop and
evaluate a bidirectional Quechua--Spanish speech translation \emph{system} for medical consultations,
chaining speech recognition (MMS), machine translation (NLLB adapted to Cusco Quechua, \sig{quz}) and
speech synthesis (MMS-TTS) through the \emph{fine-tuning} of open-source pretrained models.

The
two-way voice output aims to include Quechua-speaking people with low or no literacy. Validation will
be carried out with simulated consultations, without real patient data, measuring accuracy (WER,
BLEU, chrF), voice quality (MOS) and perceived usefulness. The study is expected to contribute the
first clinical Quechua--Spanish corpus and empirical evidence on the feasibility of these
technologies in the Andean health domain.

\vspace{0.3cm}
\noindent\obsdos{\textit{Keywords:} Machine translation, Quechua, Natural language processing,
Doctor--patient communication, Digital inclusion.}
\onehalfspacing
